% A.49 is reserved by the root for the exact prime-to-residue-characteristic comparison.
\subsection{One point has cuspidal and noncuspidal pro-prime projections}

\begin{statement}{A.49}{Refuted}{The universal prime independence in L12(ii)}
Let $s_5$ be the local section over $\Q_5$ of the rational point
$t=5$ of $U$.  In the full geometrically pro-$\ell$ arithmetic
quotient, for every $\ell\ne5$, its class is the cusp-0 section
with tangent Kummer parameter $\kappa_\ell(5)$.  In the
geometrically pro-5 quotient it belongs to no cusp packet, already
after geometric abelianization modulo 5.
Thus a single ordinary point section, ramified at 5 in the
Legendre system, does not have one cuspidal description valid for
all coefficient primes.  The universal claim in L12(ii) is false.
For an actual full cuspidal section, the projections of its
specified cusp and completed tangent parameter remain coherent
under the quotient maps, as follows directly from those maps.
\end{statement}
\begin{proof}
We prove the prime-to-residue-characteristic assertion with
explicit scheme morphisms.  Let $K/\Q_p$ be finite, let
$\mathcal O_K$ be its valuation ring with uniformizer $\pi$, and
put
\[
 R=\mathcal O_K[[T]],\qquad
 \mathcal T=\operatorname{Spec}R[1/(\pi T)].
\tag{A.49.1}
\]
There is a $K$-morphism $\mathcal T\to U_K$ from the ring map
\[
 K[t,t^{-1},(1-t)^{-1}]\longrightarrow R[1/(\pi T)],
                       \qquad t\longmapsto T.
\tag{A.49.2}
\]
It is defined because $1-T$ has the actual inverse
$\sum_{j\ge0}T^j$ in $R$.  For every nonzero
$x\in\mathfrak m_K$, evaluation of convergent integral power
series at $T=x$ gives a $K$-algebra homomorphism
$R[1/(\pi T)]\to K$.  It defines a point
$\operatorname{Spec}K\to\mathcal T$ whose composite with
(A.49.2) is precisely $t=x$.

We next compute the needed geometric group.  For $\ell\ne p$,
\[
              \pi_1(\mathcal T_{\bar K})^{(\ell)}
                           \simeq\Z_\ell(1),
\tag{A.49.3}
\]
with coordinate given by the Kummer covers $T^{1/\ell^a}$.
Here is a proof which includes the comparison behind (A.49.3).
A finite geometric Galois cover with $\ell$-group $G$ descends,
together with all its deck transformations, to a finite extension
$F/K$.  Finite etale descent over the directed union of constants
is the equivalence in Stacks, Lemma 58.14.1, tag 0BTU
\cite{LegStacksFields2026}.  Since $\mathcal O_F$ is finite free
over $\mathcal O_K$, base change gives exactly
\[
           \mathcal T_F=
           \operatorname{Spec}\mathcal O_F[[T]][1/(\pi_FT)].
\]
Normalize $\mathcal O_F[[T]]$ in the descended cover's function
field.  The normalization is finite because this complete
Noetherian local ring is excellent and Nagata
\cite[tag 07QW]{LegStacksExcellent2026}.  Away from the two
height-one primes $(\pi_F),(T)$ it is already etale.  At these
primes its ramification indices divide $|G|$, and its residue
extensions are separable: the inseparable residue degree is a
power of the residue characteristic dividing the Galois degree
\cite[tag 09EB]{LegStacksRamification2026}.  Thus it is tame
there, since $|G|$ is prime to $p$ and the residue characteristic
at $(T)$ is zero.

Choose an $\ell$-power $n$ divisible by both ramification indices,
and make the finite base change
\[
 F'=F(\mu_n,\pi_F^{1/n}),\qquad
 \mathcal O_F[[T]]\longrightarrow\mathcal O_{F'}[[u]],
                         \quad T=u^n.
\tag{A.49.4}
\]
The target is a regular complete local ring.  At both boundary
divisors the base-change ramification index is divisible by $n$;
adjoining $\mu_n$ is unramified over the local field $F$.
The discrete-valuation Abhyankar lemma says that normalized base
change is unramified when the tame ramification index divides
the new base index.  All its hypotheses have just been checked
at both height-one primes
\cite[Lemma 15.116.4, tag 0BRM]{LegStacksAbhyankar2026}.
The pulled-back normalization is therefore unramified in
codimension one and is finite over a regular ring.  Purity makes
it finite etale everywhere
\cite[Lemma 58.21.4, tag 0BMB]{LegStacksPurity2026}.

A finite etale algebra over the henselian local ring
$\mathcal O_{F'}[[u]]$ is determined by its finite etale residue
algebra.  In particular, an unramified constant extension
$F''/F'$ splits it
\cite[Lemma 10.153.7, tag 04GK]{LegStacksHenselian2026}.
Over $\bar K$ this proves that the original cover is dominated
by $u^n=T$.  Conversely, every Kummer cover
$u^{\ell^a}=T$ is geometrically connected: the $T$-valuation of
$T$ is one, so it is not an $\ell$th power after any constant
extension.  These covers have groups $\mu_{\ell^a}$ and form a
cofinal system for the geometric pro-$\ell$ group.  This proves
(A.49.3), including its cyclotomic Galois action.

We now compare this group to the actual cusp group.  The inclusion
\[
                  R[1/(\pi T)]\hookrightarrow K((T))
\tag{A.49.5}
\]
defines $\operatorname{Spec}K((T))\to\mathcal T$ over $K$.
The $T$-adic residue characteristic is zero.  Its inertia is
$\widehat\Z(1)$: after an unramified constant extension, the
same one-dimensional tame calculation gives the extensions
$T^{1/n}$.  Acting on coefficients and fixing a compatible
system of these roots gives the tangent-one splitting over
$G_K$.  Write $D_{0,\ell}$ for its quotient with inertia
$\Z_\ell(1)$.  The map (A.49.5) induces a commutative diagram
\[
\begin{array}{ccccccccc}
1&\longrightarrow&\Z_\ell(1)&\longrightarrow&D_{0,\ell}
 &\longrightarrow&G_K&\longrightarrow&1\\
&&\Vert&&\downarrow&&\Vert\\
1&\longrightarrow&\Z_\ell(1)&\longrightarrow&
 \pi_1(\mathcal T)^{\mathrm{geo},\ell}
 &\longrightarrow&G_K&\longrightarrow&1.
\end{array}
\tag{A.49.6}
\]
The geometric-and-arithmetic exact sequence here applies to
quasi-compact quasi-separated schemes; geometric connectedness
holds because each constant base change in (A.49.1) is integral
\cite[Lemma 58.14.3, tag 0BTX]{LegStacksFields2026}.
The left arrow is the identity in the Kummer coordinate $T$ by
(A.49.3).  Consequently the middle arrow is an isomorphism:
its kernel is trivial on both the geometric subgroup and the
quotient, and any element downstairs can first be lifted over
$G_K$ and then corrected by its inertia coordinate.

Composing (A.49.5) with (A.49.2) is the completion map at cusp 0
of $U_K$.  Its induced image is the cusp decomposition group,
by the valuation construction of that group on finite covers.
Its pro-$\ell$ inertia injects into the geometric group of $U_K$,
since the $t$-Kummer map is already the identity on that inertia.
Thus (A.49.6) proves that the section supplied by the point
$T=x$ maps into the actual cusp-0 packet in the full
geometrically pro-$\ell$ quotient of $\pi_1(U_K)$.
Relative to tangent one, its class in
$H^1(K,\Z_\ell(1))$ is exactly $\kappa_\ell(x)$: this follows
by pulling back the same covers $T^{1/\ell^a}$ along evaluation
at $x$.  This identifies the full section class, since the
geometric kernel of $D_{0,\ell}$ is the abelian group
$\Z_\ell(1)$.  All base-point choices change the maps by geometric
conjugation, which is precisely the convention for section classes.

Taking $K=\Q_5$, $p=5$, $x=5$ gives the asserted cuspidal
description for every $\ell\ne5$.  For $\ell=5$, the abelian
Kummer coordinates of the point section are
\[
                      (\kappa_5(5),\kappa_5(-4)).
\]
The fifth powers of units modulo 25 are $1,7,18,24$;
$-4\equiv21\pmod{25}$ is not among them.  Thus the second
coordinate is nonzero even modulo five, excluding cusp 0.
Cusp 1 would make the first coordinate zero, contradicted by
$v_5(5)=1$.  Cusp infinity would make the two coordinates equal,
since $\kappa_5(-1)=0$, contradicted by their valuations 1 and 0.
These finite checks were retained in
\texttt{code/legendre\_checks.py}.  Therefore no pro-5 cusp packet
contains the projected point section.

Finally, $E_5$ has $\Delta=6400$ of 5-adic valuation two and
$c_4=336$ a unit.  Its reduction is multiplicative of type $I_2$;
the nodal slopes are $\pm\sqrt{-1}\in\Q_5$, so it is split.
Its coefficient-$\ell$ Tate inertia is nontrivial for
$\ell\ne5$.  The counterexample is therefore at an actual place
of Legendre ramification, as required by L12.

The proved replacement is explicit: for an ordinary point in a
punctured residue disc and a coefficient prime different from the
residue characteristic, (A.49.1)--(A.49.6) identify its projection
with the corresponding cusp parameter.  The substitutions
$u=1-t$ and $u=1/t$ are actual $K$-automorphisms of $U_K$;
they give the same assertion at 1 and infinity, with parameters
$\kappa_\ell(1-x)$ and $\kappa_\ell(1/x)$ respectively.
No assertion at the coefficient prime $p$ follows from this
comparison, and the point 5 shows that the missing assertion
is false.
\end{proof}
